% !TEX root = ../Thesis.tex
\pagestyle{empty}
\hfill
\vfill
\newcommand{\coeur}[1][fill=red] % rouge par defaut
{\draw [#1] (0,0)
 .. controls +(0,2)  and +(0,2)  .. (3,0)
 .. controls +(0,-2) and +(0,2)  .. (0,-4)
 .. controls +(0,2)  and +(0,-2) .. (-3,0)
 .. controls +(0,2)  and +(0,2)  .. (0,0);}

\vfill
Nina, I liebe Di!
\centering
\begin{tikzpicture}
% mer
\shade[top color=blue,bottom color=lightgray] (0,0) arc (105:75:20)--cycle;
% quai
\fill[rotate=12] (0,0)rectangle (1,0.2);
% personnage
\fill[rotate=12,shift={(0.8,0.3)}] (0,0) ellipse (0.12 and 0.25) (0,0.35) circle (0.1);
\begin{scope} [shift={(8.4,0.4)},rotate=-10]
  % barque
  \fill (0,0)--(1,0)--(1.1,0.2)--(-0.1,0.1)--cycle;
  % rames
  \draw[line width=2pt] (-0.2,0.4)--(0.5,0.2)--(1.2,0.4);
  % personnage
  \fill[shift={(0.5,0.25)},rotate=-10]
     (0,0)ellipse(0.12 and 0.25) (0,0.35)circle(0.1);
\end{scope}
% coeurs
\coeur[shift={(1.3,2)},ball color=red,scale=0.35]
\coeur[shift={(0.3,1.8)},ball color=blue,scale=0.2,rotate=30*rand]
\coeur[shift={(2.2,1.6)},ball color=magenta,scale=0.2,rotate=30*rand]
\coeur[shift={(1.8,2.5)},ball color=green,scale=0.25,rotate=30*rand]
\coeur[shift={(2.7,2.1)},ball color=yellow,scale=0.15,rotate=30*rand]
\end{tikzpicture}
\vfill